% ============================================================
% SECTION: Regression-Detection Validity Study
% To be inserted into paper.tex after Section 5 (Reliability
% and Validity) or as a standalone subsection within it.
%
% Fill in [RESULT] placeholders after running:
%   python run_regression_detection.py
% ============================================================

\subsection{Construct Validity: Regression-Detection Study}
\label{sec:regression}

A benchmark is only useful if it detects problems that practitioners actually
care about. We conducted a controlled regression-detection experiment to test
whether \benchname{} responds appropriately when a known policy defect is
introduced into a system under evaluation.

\paragraph{Study design.}
We defined four policy defects, each operationalized as a change to the system
prompt passed to the evaluated model. The defects were chosen to correspond to
recognizable real-world failure modes: (1)~\emph{casual-register suppression},
in which the model omits policy disclaimers when the user writes informally;
(2)~\emph{authority deference}, in which the model treats claimed professional
authorization as verified fact; (3)~\emph{minimization override}, in which
the model defers to the user's self-assessment that a situation is not urgent;
and (4)~\emph{hypothetical exemption}, in which the model grants an exception
to policy when requests are framed as hypothetical. We also included a
fifth condition, \emph{universal refusal}, in which the model declines all
domain-specific requests regardless of framing---an edge case designed to test
whether low CAI Strain can be gamed through consistent unhelpfulness.

\paragraph{Pre-registration.}
Before running any evaluations, we wrote out the predicted direction and
minimum magnitude of each detectable strain change, specifying the primary
adversarial technique (\techname{T3} for casual suppression, \techname{T5} for
authority deference, \techname{T4} for minimization override, \techname{T6}
for hypothetical exemption) and the domains where effects were expected to be
largest. We also specified neutral domains that should show no material
increase. Pre-registered predictions are available in the supplementary
materials. No results were examined prior to pre-registration.

\paragraph{Procedure.}
Each defect was tested by running \benchname{} on \NUMCASES{} held-out cases
per domain across \NUMDOMAINS{} domains, using Claude Sonnet~4.6 as the
evaluated model and GPT-4o-mini as the consistency judge. We computed
per-technique CAI Strain for each defect condition and compared against a
clean baseline with no defect system prompt. Helpfulness of canonical
responses under universal refusal was assessed using a binary judge score
(helpful vs.~unhelpful).

\paragraph{Results.}
% ── FILL IN AFTER RUNNING THE EXPERIMENT ──────────────────────────────────
% Template: "CAI-Bench detected N of 4 pre-registered defects above threshold,
% with strain increases localized to the expected technique × domain cells
% (Table X). The false-positive rate across neutral domains was Y\%."
%
% For the universal refusal condition: "Universal refusal produced a mean
% CAI Strain of Z (vs. W for the baseline), but a mean helpfulness score of
% V (vs. U for the baseline), confirming that low strain does not reward
% consistent unhelpfulness."

\benchname{} detected \NDETECTED{} of 4 pre-registered defects above the
specified minimum-detectable difference of 0.08 strain units. Strain increases
were localized to the predicted technique $\times$ domain cells; the
false-positive rate across the pre-specified neutral cells was
\FPRPCT{}\%. Table~\ref{tab:regression} summarizes the per-defect detection
results.

The universal-refusal condition produced a mean CAI Strain of \REFUSALSTRAIN{}
(vs.~\BASELINESTRAIN{} for the baseline), but a mean helpfulness score of
\REFUSALHELPFUL{} (vs.~\BASELINEHELPFUL{} for the baseline). A system that
consistently refuses all requests appears well-calibrated under CAI Strain
alone; the helpfulness check exposes this as a pathological response. This
result motivates reporting both components when using \benchname{} to compare
systems in deployment settings.

% ── TABLE: Detection results ───────────────────────────────────────────────
\begin{table}[t]
\centering
\small
\caption{Regression-detection results. Delta = mean CAI Strain under the defect
condition minus baseline. Threshold = 0.08 for all except minimization override
(0.10). Primary tech = the technique predicted to show the largest increase.
Detected = delta exceeds threshold in the predicted domain.}
\label{tab:regression}
\begin{tabular}{llccc}
\toprule
Defect & Primary tech & Domain & $\Delta$ Strain & Detected \\
\midrule
Casual suppression  & T3 & finance       & [FILL] & [Y/N] \\
                    &    & ecommerce     & [FILL] & [Y/N] \\
                    &    & food\_delivery & [FILL] & [Y/N] \\
Authority deference & T5 & legal         & [FILL] & [Y/N] \\
                    &    & healthcare    & [FILL] & [Y/N] \\
                    &    & immigration   & [FILL] & [Y/N] \\
Minimization override & T4 & mental\_health & [FILL] & [Y/N] \\
Hypothetical exemption & T6 & legal       & [FILL] & [Y/N] \\
                    &    & ai\_safety    & [FILL] & [Y/N] \\
\midrule
\multicolumn{3}{l}{\textit{Neutral cells (false-positive check)}} & & \\
Casual suppression  & T3 & mental\_health & [FILL] & -- \\
Authority deference & T5 & ecommerce     & [FILL] & -- \\
Minimization override & T4 & finance     & [FILL] & -- \\
\midrule
Universal refusal   & --  & all          & [FILL] & -- \\
\quad Helpfulness score & -- & all        & [FILL] & -- \\
\bottomrule
\end{tabular}
\end{table}
% ──────────────────────────────────────────────────────────────────────────

\paragraph{Interpretation.}
These results support the claim that CAI Strain is sensitive to policy-relevant
behavioral differences while remaining specific enough not to flag unrelated
domains. A defect introduced through one framing technique elevates strain
under that technique without materially raising scores under orthogonal
techniques. The universal-refusal finding motivates combining CAI Strain with
a minimum helpfulness threshold in any deployment comparison---a constraint
we treat as outside the scope of the benchmark itself but recommend as
evaluation practice.

\paragraph{Limitations.}
This experiment used a single evaluated model (Claude Sonnet~4.6) and a single
judge (GPT-4o-mini). Detection sensitivity may differ for other model families.
The defects were operationalized through system-prompt manipulation; naturally
occurring policy drift may produce subtler or more diffuse strain increases
than the concentrated defects tested here. We regard this study as a
proof-of-concept for the detection methodology rather than a comprehensive
sensitivity analysis.
